\chapter{Barisan dan Deret}\label{ch:b2:series}

Teori deret numerik (jilid Tahun ke-1) matang di sini pada tiga
arah: deret yang bernilai di \omterm{def:b2:nvs:banach}{ruang Banach}, tempat kekonvergenan
\omterm{def:b2:series:def}{mutlak} melakukan pekerjaannya; uji yang lebih halus bagi deret real
(\omterm{thm:b2:series:abel}{penjumlahan Abel}); serta \emph{\omterm{def:b2:series:summable}{keluarga terjumlahkan}} --- penjumlahan
yang dibebaskan dari urutan sukunya --- \omterm{def:b2:metric:complete}{lengkap} dengan teorema Fubini
untuk jumlah ganda dan hasil kali Cauchy. Perkakas ini mengusung semua
bab deret fungsi berikutnya.

\section{Deret pada ruang bernorma}

\begin{definition}\label{def:b2:series:def}
Untuk barisan $(u_n)$ pada ruang bernorma $E$, deret $\sum u_n$
konvergen apabila jumlah parsialnya konvergen; deret itu konvergen
\emph{mutlak}\index{kekonvergenan mutlak!pada ruang Banach} apabila
$\sum \norm{u_n} < \infty$. Pada ruang \emph{Banach}, kekonvergenan
mutlak mengakibatkan kekonvergenan
(\cref{thm:b2:nvs:absoluteconvergence}); sedangkan pada ruang yang tidak
\omterm{def:b2:metric:complete}{lengkap} hal itu dapat gagal (\cref{exo:b2:series:9}).
\end{definition}

\begin{example}\label{ex:b2:series:neumann}
Di $\mathcal{M}_n(K)$ (atau $\mathcal{L}_c(E)$ dengan $E$ Banach): untuk
$\vertiii A < 1$, \emph{deret Neumann} $\sum A^k$ konvergen
\omterm{def:b2:series:def}{mutlak} ke $(I - A)^{-1}$ (yang dibuktikan pada \cref{exo:b2:nvs:5});
sedangkan $\sum \frac{A^k}{k!}$ konvergen \omterm{def:b2:series:def}{mutlak} ke $\eu^A$ untuk setiap
$A$ (\cref{ex:b2:nvs:matrixexp}). Deret geometri dan deret
eksponensial bernilai operator berkelakuan seperti model skalarnya ---
dan itulah seluruh inti kerangka Banach.
\end{example}

\section{Penjumlahan Abel}

\begin{theorem}[Penjumlahan Abel dan ujinya]\label{thm:b2:series:abel}
(Penjumlahan parsial) Untuk skalar $a_n$ dan vektor $b_n$, dengan $B_n
= \sum_{k=0}^{n} b_k$:\index{penjumlahan Abel}
\[
\sum_{n=0}^{N} a_n b_n
= a_N B_N - \sum_{n=0}^{N-1} (a_{n+1} - a_n) B_n .
\]
(Uji Abel) Jika $(a_n)$ barisan real yang turun ke $0$, dan
jumlah parsial $B_n$ bersifat \emph{terbatas} (di sebuah \omterm{def:b2:nvs:banach}{ruang Banach}),
maka $\sum a_n b_n$ konvergen.
\end{theorem}

\begin{proof}
Kesamaannya, langkah demi langkah: dengan $B_{-1} = 0$, tulislah $b_n =
B_n - B_{n-1}$ lalu pecah,
\[
\sum_{n=0}^{N} a_nb_n
= \sum_{n=0}^{N} a_nB_n - \sum_{n=0}^{N} a_nB_{n-1}
= \sum_{n=0}^{N} a_nB_n - \sum_{n=0}^{N-1} a_{n+1}B_{n} ,
\]
dengan jumlah keduanya diindeks ulang lewat $n \mapsto n + 1$ (dan suku
$B_{-1}$ lenyap); setelah jangkauan bersamanya $0 \leq n \leq N-1$
dikumpulkan, tersisa $a_NB_N$ ditambah $\sum_{n\leq N-1}(a_n -
a_{n+1})B_n$: itulah rumus yang dinyatakan tadi. Ia adalah pengintegralan
parsial versi diskret, dengan $(B_n)$ sebagai antiturunan $(b_n)$ dan
selisih $a_{n+1} - a_n$ sebagai turunan $(a_n)$. Untuk ujinya, dengan
$\norm{B_n} \leq M$: suku batasnya $a_N B_N \to 0$; sedangkan deret
$\sum (a_n - a_{n+1})B_n$ konvergen \omterm{def:b2:series:def}{mutlak}, sebab
\[
\sum_n \norm{(a_{n+1} - a_n)B_n} \leq M \sum_n (a_n - a_{n+1})
= M a_0 < \infty
\]
(lewat teleskop, dengan $a_n \downarrow 0$). Kedua potongan kesamaannya
konvergen, sehingga $\sum a_n b_n$ pun konvergen.
\end{proof}

\begin{example}\label{ex:b2:series:sinn}
Deret $\sum \frac{\sin n}{n}$ konvergen: sebab $a_n = \frac1n
\downarrow 0$ dan $B_n = \sum_{k=1}^{n} \sin k$ terbatas --- memang $B_n
= \Im\sum_{k \leq n} \eu^{\iu k} = \Im\,\frac{\eu^{\iu}(\eu^{\iu n} -
1)}{\eu^{\iu} - 1}$, yang bermodulus $\leq \frac{2}{\abs{\eu^{\iu} -
1}}$. Deret itu \emph{tidak} konvergen \omterm{def:b2:series:def}{mutlak} (sebab $\abs{\sin n} \geq
\sin^2 n = \frac{1 - \cos 2n}{2}$, dan $\sum \frac{1 - \cos 2n}{2n}$
divergen karena $\sum \frac{\cos 2n}{n}$ konvergen lewat uji Abel yang
sama sedangkan $\sum \frac{1}{2n}$ divergen). Uji deret
berselang-seling tak lain kasus khusus $b_n = (-1)^n$.
\end{example}

\begin{example}[Abel pada lingkaran kekonvergenan]\label{ex:b2:series:abelboundary}
Untuk $z$ kompleks yang mana dengan $\abs z = 1$ deret $\sum_{n \geq 1}
\frac{z^n}{n}$ konvergen? Di $z = 1$ ia menjadi deret harmonik, jadi
divergen. Untuk $z \neq 1$ pada lingkarannya, uji Abel berlaku
dengan $a_n = \frac1n \downarrow 0$ dan $b_n = z^n$, yang jumlah
parsialnya terbatas tanpa bergantung pada $N$:
\[
\Bigl|\sum_{n=1}^{N} z^n\Bigr|
= \Bigl|\frac{z(z^N - 1)}{z - 1}\Bigr|
\leq \frac{2}{\abs{z - 1}} .
\]
Jadi konvergen --- walaupun tak pernah \omterm{def:b2:series:def}{mutlak} (sebab $\sum\frac1n$).
Satu deret, selingkar perilaku: divergen di satu titik saja, dan
semi-konvergen di semua titik lainnya. Inilah perilaku baku di
perbatasan bagi deret pangkat (\cref{ch:b2:powerseries}), yang di sini
dijumpai dengan tangan kosong; di $z = -1$ ia memulihkan deret
harmonik berselang-seling, dan di $z = \eu^{\iu\theta}$ bagian real dan
bagian imajinernya adalah deret $\sum\frac{\cos
n\theta}{n}$ dan $\sum\frac{\sin n\theta}{n}$ pada
\cref{exo:b2:series:4}.
\end{example}

\begin{example}[Deret berselang-seling yang beranjau]\label{ex:b2:series:trap}
Apakah $\sum_{n\geq2} \dfrac{(-1)^n}{\sqrt n + (-1)^n}$ konvergen?
Tandanya berselang-seling dan sukunya menuju $0$ --- namun
uji berselang-selingnya \emph{tidak} berlaku: sebab modulusnya
$\frac{1}{\sqrt n + (-1)^n}$ tidak turun (ia melonjak naik pada
tiap $n$ yang ganjil). Uraikanlah sebagai gantinya:
\[
\frac{(-1)^n}{\sqrt n + (-1)^n}
= \frac{(-1)^n}{\sqrt n}\cdot
\frac{1}{1 + \frac{(-1)^n}{\sqrt n}}
= \frac{(-1)^n}{\sqrt n} - \frac{1}{n}
+ O\Bigl(\frac{1}{n^{3/2}}\Bigr).
\]
Potongan pertamanya konvergen (lewat uji berselang-seling yang
diterapkan secara jujur pada $\frac1{\sqrt n}\downarrow0$), potongan
ketiganya konvergen \omterm{def:b2:series:def}{mutlak} --- tetapi potongan tengahnya adalah deret
harmonik yang divergen: jadi jumlahnya \emph{divergen} ke $-\infty$.
Pelajaran penutupnya: ketika kemonotonannya gagal, uraikanlah sampai tiap
potongannya konvergen \omterm{def:b2:series:def}{mutlak} atau menjadi kasus uji yang bersih; sebab
$-\frac1n$ yang tersembunyi itu tak terlihat oleh pencacahan tanda.
\end{example}

\begin{remark}[Jebakan yang sering muncul]\label{rem:b2:series:pitfalls}
(i) ``Sukunya menuju $0$'' tidak membuktikan apa pun: sebab deret
harmonik divergen. (ii) Uji berselang-seling menuntut modulus yang
\emph{turun} --- \cref{ex:b2:series:trap} adalah contoh penyangkal
kanoniknya, dan deret ketiga pada \cref{exo:b2:series:1} menjadi
latihannya. (iii) Deret yang konvergen bersyarat tidak boleh
ditata ulang (\cref{ex:b2:series:rearrange}), dan hasil kali
Cauchy-nya dapat divergen: untuk $\sum\frac{(-1)^n}{\sqrt{n+1}}$
yang dikuadratkan, suku diagonalnya memenuhi
\[
\abs{c_k} = \sum_{m=0}^{k}
\frac{1}{\sqrt{(m+1)(k-m+1)}}
\geq (k+1)\cdot\frac{2}{k+2} \longrightarrow 2 \neq 0
\]
(sebab tiap faktornya paling banyak $\frac{k+2}2$ menurut AM--GM),
sehingga $\sum c_k$ divergen --- jadi kekonvergenan \omterm{def:b2:series:def}{mutlak} sekurangnya
satu faktor (\cref{exo:b2:series:8}) bukanlah kemewahan. (iv)
Keterjumlahan menurut definisinya berbicara tentang batas yang
\emph{\omterm{def:b2:series:def}{mutlak}}: jadi tak ada yang namanya \omterm{def:b2:series:summable}{keluarga terjumlahkan} bersyarat.
\end{remark}

\section{Keluarga terjumlahkan}

\begin{definition}\label{def:b2:series:summable}
Misalkan $I$ himpunan indeks yang \omterm{def:b2:structures:countable}{terbilang}. Sebuah keluarga $(u_i)_{i
\in I}$ berisi \emph{bilangan real tak negatif} disebut
\emph{terjumlahkan}\index{keluarga terjumlahkan} apabila jumlah parsial
hingganya terbatas; jumlahnya adalah
\[
\sum_{i \in I} u_i = \sup_{F \subseteq I \text{ hingga}} \sum_{i
\in F} u_i \in \intcc{0}{+\infty} .
\]
Sebuah keluarga bilangan real atau kompleks (atau vektor Banach) disebut
terjumlahkan apabila $(\norm{u_i})$ demikian; jumlahnya lalu ditetapkan
dengan memecahnya menjadi bagian positif dan negatif (atau bagian real
dan imajiner) --- setara dengan itu, sebagai nilai bersama $\sum_{n}
u_{\sigma(n)}$ atas semua pencacahan $\sigma$ bagi $I$ (lihat di bawah).
\end{definition}

\begin{method}[Memilih sebuah uji]\label{met:b2:series:choosing}
Menghadapi $\sum u_n$, urutannya begini. (1) Jika $u_n \not\to 0$, ia
divergen, berhenti. (2) Jika sukunya bertanda tetap, bandingkanlah:
carilah yang setara (\cref{ch:b2:comparison}) lalu tempatkan ia pada
peta Riemann--Bertrand. (3) Jika tandanya berselang-seling dengan
modulus yang \emph{turun}, pakai uji berselang-seling; sedangkan jika
modulusnya tidak monoton, uraikanlah sukunya sampai tiap potongannya
konvergen \omterm{def:b2:series:def}{mutlak} atau menjadi kasus uji yang bersih
(\cref{ex:b2:series:trap}). (4) Jika pola tandanya berayun tetapi
terstruktur ($\sin n\theta$, $\eu^{\iu n\theta}$, pangkat matriks),
pakai uji Abel dengan jumlah parsial yang terbatas. (5) Kekonvergenan
\omterm{def:b2:series:def}{mutlak} selalu patut diperiksa lebih dulu: ia lebih kuat, kebal terhadap
urutan, dan membuka hasil kali Cauchy beserta Fubini.
\end{method}

\begin{example}[Keterjumlahan lewat pencacahan diagonal]\label{ex:b2:series:diagonalcount}
Untuk $s > 0$ yang mana keluarga
$\bigl((m + n)^{-s}\bigr)_{m, n \geq 1}$ \omterm{def:b2:series:summable}{terjumlahkan}? Kelompokkan
jumlah parsial hingganya menurut diagonal $m + n = k$: diagonal ke-$k$
mengusung $k - 1$ pasangan, yang masing-masing menyumbang $k^{-s}$, jadi
jumlah hingganya tepat terbatas oleh (dan menghabiskan)
\[
\sum_{k \geq 2} \frac{k - 1}{k^{s}} ,
\]
yakni deret bersuku positif yang setara dengan $k^{1-s}$: jadi
\omterm{def:b2:series:summable}{terjumlahkan} bila dan hanya bila $s - 1 > 1$, yakni $s > 2$. Indeks
berdimensi dua memakan satu pangkat penuh: sebuah bidang suku ``lebih
divergen satu dimensi'' daripada sebuah garis --- jadi geometri
pencacahan himpunan indeksnya, bukan besarnya tiap suku, yang
memutuskan keterjumlahannya.
(Pencacahan yang sama menunjukkan $\bigl((m^2 + n^2)^{-1}\bigr)$ tidak
\omterm{def:b2:series:summable}{terjumlahkan}: pada diagonal $m + n = k$, tiap sukunya sedikitnya
$k^{-2}$, dan $(k-1)\cdot k^{-2}$ berjumlah seperti deret
harmonik.)
\end{example}

\begin{theorem}[Keterjumlahan dan urutan]\label{thm:b2:series:rearrangement}
\begin{enumerate}
  \item Untuk keluarga tak negatif, jumlahnya awet terhadap
        pencacahan apa pun: $\sum_{i} u_i = \sum_{n=0}^{\infty}
        u_{\sigma(n)}$ untuk setiap bijeksi $\sigma \colon \N \to I$.
  \item Deret real atau kompleks $\sum u_n$ disebut \emph{konvergen
        komutatif} (yakni setiap penataan ulangnya konvergen, dengan
        jumlah yang sama) bila dan hanya bila ia konvergen \omterm{def:b2:series:def}{mutlak}.
\end{enumerate}
\end{theorem}

\begin{proof}
(1) Setiap jumlah parsial $\sum_{n \leq N} u_{\sigma(n)}$ merupakan
jumlah parsial hingga keluarganya (jadi $\leq$ supremumnya); sebaliknya
setiap $F$ yang hingga termuat di suatu $\{\sigma(0), \dots,
\sigma(N)\}$ (jadi supremumnya $\leq$ limit deretnya). Kedua batasnya berimpit.

(2) Jika $\sum\abs{u_n} < \infty$: untuk sembarang penataan ulang
$\sigma$ dan $\varepsilon > 0$, pilihlah $N$ dengan $\sum_{n >
N}\abs{u_n} \leq \varepsilon$; maka melampaui peringkat tempat $\sigma$
sudah menghabiskan $\intint{0}{N}$, jumlah parsial yang ditata ulang itu
berselisih dari limit asalnya paling banyak $\varepsilon$: jadi jumlahnya
sama. Jika $\sum \abs{u_n} = \infty$ padahal $\sum u_n$ konvergen (kasus
real; kasus kompleksnya menyusul koordinat demi koordinat): maka bagian
positif dan bagian negatifnya sama-sama divergen, dan penataan ulangnya
dapat mencapai limit apa pun yang ditentukan --- itulah teorema Riemann,
yang dikerjakan pada \cref{exo:b2:series:5} --- sehingga kekonvergenan
komutatifnya gagal.
\end{proof}

\begin{example}[Sebuah penataan ulang tertangkap basah]\label{ex:b2:series:rearrange}
Deret harmonik berselang-seling berjumlah
$\sum_{n\geq1}\frac{(-1)^{n-1}}{n} = \ln 2$ (jilid Tahun
ke-1). Tatalah ulang sebagai ``satu positif, dua negatif'':
\[
1 - \frac12 - \frac14 + \frac13 - \frac16 - \frac18 + \frac15 -
\cdots
\]
Setelah tiap blok bertiganya dikelompokkan,
\[
\frac{1}{2k-1} - \frac{1}{4k-2} - \frac1{4k}
= \frac{1}{4k-2} - \frac{1}{4k}
= \frac12\Bigl(\frac{1}{2k-1} - \frac1{2k}\Bigr),
\]
maka deret yang ditata ulang itu konvergen ke $\frac12\ln 2$ --- separuh
jumlah asalnya, dengan suku yang persis sama. Deret yang tidak
konvergen \omterm{def:b2:series:def}{mutlak} mengingat urutan sukunya; sedangkan \omterm{def:b2:series:summable}{keluarga
terjumlahkan} justru yang tidak mengingatnya.
\end{example}

\begin{example}[Mengelompokkan itu aman, membuka kelompok tidak]\label{ex:b2:series:grouping}
Mengelompokkan suku berurutan sebuah deret yang \emph{konvergen} tak
pernah mengubah jumlahnya: sebab jumlah parsial yang dikelompokkan
membentuk barisan bagian jumlah parsial asalnya. Operasi sebaliknya
terlarang:
\[
(1 - 1) + (1 - 1) + (1 - 1) + \cdots = 0 + 0 + \cdots = 0,
\]
padahal $1 - 1 + 1 - 1 + \cdots$ yang tanpa kelompok bersifat divergen
(sebab jumlah parsialnya berayun antara $1$ dan $0$). Membuka kelompok
hanya sah bila ada hipotesis pengimbang --- misalnya suku yang
menuju $0$ dengan panjang kelompok yang terbatas: sebab ketika itu, di
antara dua jumlah parsial yang dikelompokkan, jumlah asalnya hanyut
paling banyak sejumlah suku $o(1)$ yang banyaknya terbatas, sehingga
kekonvergenannya berpindah kembali. Persis itulah klausul yang membuat
perhitungan blok pada \cref{ex:b2:series:rearrange} menjadi sebuah bukti,
bukan sulap.
\end{example}

\begin{theorem}[Fubini untuk keluarga; hasil kali Cauchy]\label{thm:b2:series:fubini}
Misalkan $(u_{m,n})_{(m,n) \in \N^2}$ keluarga ganda yang \omterm{def:b2:series:summable}{terjumlahkan}
(yakni $\sup_F \sum_F \abs{u_{m,n}} < \infty$). Maka\index{teorema
Fubini untuk deret}
\[
\sum_{(m,n)} u_{m,n}
= \sum_{m=0}^{\infty}\Bigl(\sum_{n=0}^{\infty} u_{m,n}\Bigr)
= \sum_{n=0}^{\infty}\Bigl(\sum_{m=0}^{\infty} u_{m,n}\Bigr)
= \sum_{k=0}^{\infty}\Bigl(\sum_{m+n=k} u_{m,n}\Bigr),
\]
dengan semua deret dalamnya konvergen (secara \omterm{def:b2:series:def}{mutlak}). Khususnya, jika
$\sum a_m$ dan $\sum b_n$ konvergen \omterm{def:b2:series:def}{mutlak}, maka \emph{hasil kali
Cauchy}\index{hasil kali Cauchy} keduanya konvergen \omterm{def:b2:series:def}{mutlak} dengan
\[
\Bigl(\sum_m a_m\Bigr)\Bigl(\sum_n b_n\Bigr)
= \sum_{k=0}^{\infty} c_k,
\qquad
c_k = \sum_{m=0}^{k} a_m b_{k-m} .
\]
\end{theorem}

\begin{proof}
\emph{Kasus tak negatif.} Tiap pengelompokan (menurut baris, kolom, atau
diagonal) menghitung supremum yang sama: sebab sembarang himpunan
pasangan yang hingga termuat di sebuah blok baris yang hingga (sehingga
tiap jumlah kelompoknya terbatas di bawah oleh jumlah parsial hingga dan
di atas oleh totalnya), lalu kekonvergenan monoton jumlah parsialnya
mengerjakan sisanya --- secara konkret, untuk barisnya: dari $\sum_{m
\leq M}\sum_{n \leq N} u_{m,n} \leq S$, dengan melewatkan $N \to \infty$
lalu $M \to \infty$, diperoleh $\sum_m \sum_n u_{m,n} \leq S$;
sebaliknya setiap $F$ yang hingga duduk di dalam persegi panjang
semacam itu, jadi $S \leq \sum_m\sum_n u_{m,n}$. Untuk diagonalnya:
kedua batas yang sama, dengan segitiga menggantikan persegi panjangnya.

\emph{Kasus umumnya.} Pecahlah menjadi bagian positif dan negatif (atau
real dan imajiner), yang masing-masing keluarga tak negatif yang
\omterm{def:b2:series:summable}{terjumlahkan}; keempat pengelompokannya sepakat pada tiap bagiannya, jadi
sepakat pula pada selisihnya; sedangkan kekonvergenan \omterm{def:b2:series:def}{mutlak} deret
dalamnya berasal dari kasus tak negatif yang diterapkan pada
$\abs{u_{m,n}}$.

\emph{\omterm{thm:b2:series:fubini}{Hasil kali Cauchy}.} Keluarga $u_{m,n} = a_m b_n$ \omterm{def:b2:series:summable}{terjumlahkan}:
sebab jumlah parsial hingga $\abs{a_mb_n}$ terbatas oleh
$\bigl(\sum\abs{a_m}\bigr)\bigl(\sum\abs{b_n}\bigr)$. Barisnya memberi
$\bigl(\sum a_m\bigr)\bigl(\sum b_n\bigr)$; sedangkan diagonalnya
memberi $\sum_k c_k$.
\end{proof}

\begin{example}[Kesamaan eksponensial, secara jujur]\label{ex:b2:series:exp}
Untuk $a, b \in \C$ (atau matriks yang komutatif):
\[
\Bigl(\sum_m \frac{a^m}{m!}\Bigr)\Bigl(\sum_n
\frac{b^n}{n!}\Bigr)
= \sum_k \sum_{m+n=k} \frac{a^m b^n}{m!\,n!}
= \sum_k \frac{(a + b)^k}{k!},
\]
lewat teorema binomial pada tiap diagonalnya: jadi $\eu^a \eu^b =
\eu^{a+b}$ --- yakni persamaan fungsional $\exp$ yang diturunkan dari
deretnya semata. (Kekomutatifannya dipakai pada langkah binomialnya;
untuk matriks yang tidak komutatif kesamaan itu sungguh gagal,
\cref{ch:b2:diffeq}.)
\end{example}

\begin{example}[Hasil kali Cauchy sebagai alat hitung]\label{ex:b2:series:squaresum}
Dari deret geometri dan $\sum_{n \geq 1} nz^n = \frac{z}{(1-z)^2}$
($\abs z < 1$) pada \cref{exo:b2:series:2}, satu hasil kali
Cauchy lagi merampungkan momen keduanya. Kalikan $\sum_m
mz^m$ dengan $\sum_n z^n$: koefisien diagonalnya adalah $\sum_{m=0}^k
m = \frac{k(k+1)}2$, sehingga
\[
\frac{z}{(1-z)^3} = \sum_{k\geq0}\frac{k(k+1)}{2}\,z^k ,
\]
lalu kesamaan $n^2 = 2\cdot\frac{n(n+1)}2 - n$ merakit
\[
\sum_{n\geq1} n^2z^n = \frac{2z}{(1-z)^3} - \frac{z}{(1-z)^2}
= \frac{z(1+z)}{(1-z)^3} .
\]
Di $z = \frac12$: $\sum_{n\geq1}\frac{n^2}{2^n} =
\frac{\frac12\cdot\frac32}{\frac18} = 6$ --- sebuah nilai tertutup tanpa
penurunan di mana pun, hanya deret yang konvergen \omterm{def:b2:series:def}{mutlak} lalu
dikalikan seperti polinomial. Rangkaian kesamaan yang sama
menghitung setiap $\sum n^dz^n$, dan para ahli peluang akan mengenali
momen faktorial kedua pada distribusi geometri
(\cref{ch:b2:genfun}).
\end{example}

\begin{example}[Sebuah penilaian jumlah ganda]\label{ex:b2:series:doublesum}
Untuk $s > 1$ yang real, misalkan $\zeta(s) = \sum_{n\geq1} n^{-s}$.
Dengan mencacah pembagi lewat penjumlahan ganda --- keluarga
$(m^{-s}n^{-s})$ atas $(m,n) \in (\N^*)^2$ \omterm{def:b2:series:summable}{terjumlahkan} (sebab hasil
kali deret positif yang konvergen) --- lalu mengelompokkannya menurut
hasil kali $q = mn$:
\[
\zeta(s)^2 = \sum_{m,n} \frac{1}{(mn)^s}
= \sum_{q=1}^{\infty} \frac{d(q)}{q^s},
\]
dengan $d(q)$ menyatakan banyaknya pembagi $q$. \omterm{def:b2:series:summable}{Keluarga terjumlahkan}
mengubah kombinatorika menjadi analisis.
\end{example}

\begin{example}[Sebuah penilaian Fubini: $\sum_n (\zeta(n) - 1) = 1$]\label{ex:b2:series:zetasum}
Untuk bilangan bulat $n \geq 2$ berlaku $\zeta(n) - 1 = \sum_{k \geq 2}
k^{-n}$. Keluarga ganda $(k^{-n})_{k, n \geq 2}$ \omterm{def:b2:series:summable}{terjumlahkan}: dengan
menjumlahkan kolom geometrinya lebih dulu,
\[
\sum_{k\geq2}\sum_{n\geq2} \frac{1}{k^n}
= \sum_{k\geq2} \frac{1/k^2}{1 - 1/k}
= \sum_{k\geq2} \frac{1}{k(k-1)} = 1
\]
(lewat teleskop), dan semua sukunya positif, jadi
\cref{thm:b2:series:fubini} mengizinkan penjumlahan menurut barisnya:
\[
\sum_{n\geq2}\bigl(\zeta(n) - 1\bigr) = 1 .
\]
Nilai $\zeta$ yang tak hingga banyaknya, yang masing-masing tampak
transenden, punya ekor yang berjumlah tepat $1$. Pelajaran penutupnya:
ketika sebuah jumlah ganda bersuku positif, hitunglah ia dalam urutan
mana pun yang runtuh --- di sini kolomnya geometri, barisnya misterius,
dan Fubini memindahkan keruntuhannya.
\end{example}

\begin{example}[Deret geometri menyelesaikan sebuah persamaan]\label{ex:b2:series:neumannsolve}
Di \omterm{def:b2:nvs:banach}{ruang Banach} $\bigl(C(\intcc01),
\norm\cdot_\infty\bigr)$, selesaikan $x - K(x) = y$ dengan $K(f)$
menyatakan fungsi konstan $\frac12\int_0^1 f$. \omterm{thm:b2:nvs:continuouslinear}{Norma operatornya}
$\vertiii K \leq \frac12 < 1$, jadi deret Neumann berlaku
(\cref{ex:b2:series:neumann}): $x = \sum_{n\geq0} K^n(y)$.
Hitung iterasinya: $K(y) = \frac12\int_0^1 y$ (sebuah konstanta),
dan menerapkan $K$ pada konstanta $c$ memberi $\frac c2$, sehingga
$K^n(y) = \frac{1}{2^{n-1}}\cdot\frac12\int_0^1 y$ untuk $n \geq
1$. Setelah konstanta geometrinya dijumlahkan:
\[
x = y + \Bigl(\int_0^1 y\Bigr)
\sum_{n\geq1}\frac{1}{2^n} = y + \int_0^1 y .
\]
Periksa: $x - K(x) = y + \int y - \frac12\bigl(\int y + \int
y\bigr) = y$. Sebuah deret tak hingga, sebuah jawaban hingga, dan
pemeriksaan satu baris --- jadi deret geometri adalah algoritme
pembalikan, bukan sekadar pernyataan kekonvergenan.
\end{example}

\begin{example}[Teleskop lewat pecahan parsial]\label{ex:b2:series:telescope}
Penjumlahan eksak itu langka; teleskop adalah pemasok utamanya.
Uraikanlah
\[
\frac{1}{n(n+1)(n+2)}
= \frac{1}{2}\Bigl(\frac{1}{n(n+1)} -
\frac{1}{(n+1)(n+2)}\Bigr),
\]
(periksa dengan menyamakan penyebutnya), sehingga jumlah
parsialnya runtuh:
\[
\sum_{n=1}^{N}\frac{1}{n(n+1)(n+2)}
= \frac12\Bigl(\frac{1}{1\cdot2} -
\frac{1}{(N+1)(N+2)}\Bigr)
\longrightarrow \frac14 .
\]
Pola yang sama --- yakni menulis sukunya sebagai $c(u_n - u_{n+1})$
untuk suatu $(u_n)$ yang gamblang --- menyelesaikan \cref{exo:b2:series:10}
(lewat arkus tangen) dan menghitung setiap $\sum\frac{1}{n(n+1)\cdots(n +
k)} = \frac{1}{k\cdot k!}$. Bila ada jumlah eksak pada
tingkat ini, biasanya ada teleskop yang bersembunyi di dalam sukunya.
\end{example}

\begin{remark}[Pandangan ke depan di dalam jilid ini]\label{rem:b2:series:perspectives}
Tiga bab berikutnya menjadi pelanggan langsungnya. Untuk
\cref{ch:b2:funcseq}: kekonvergenan normal $\sum f_n$ tak lain
kekonvergenan \omterm{def:b2:series:def}{mutlak} $\sum\norm{f_n}_\infty$ di \omterm{def:b2:nvs:banach}{ruang Banach}
$\bigl(C, \norm\cdot_\infty\bigr)$ --- yakni
\cref{thm:b2:nvs:absoluteconvergence} yang berkostum. Untuk
\cref{ch:b2:powerseries}: di dalam cakram kekonvergenannya
segalanya \omterm{def:b2:series:def}{mutlak} dan \omterm{def:b2:series:summable}{terjumlahkan}, jadi hasil kali Cauchy dan penataan
ulang berjalan bebas (dan itulah sebabnya deret pangkat berkali seperti
polinomial); sedangkan di perbatasannya uji Abel mengambil alih
(\cref{ex:b2:series:abelboundary}). Untuk \cref{ch:b2:genfun}:
fungsi pembangkit peluang adalah deret pangkat yang segala
manipulasinya --- hasil kali untuk jumlah peubah saling bebas, jumlah
ganda untuk distribusi majemuk --- disahkan oleh
\cref{thm:b2:series:fubini}. \omterm{def:b2:series:summable}{Keluarga terjumlahkan} adalah bagian
hukum bagi analisis yang akan datang.
\end{remark}

\begin{remark}[Di mana bab ini dipakai]
Segala yang punya jumlah tak hingga melewati bab ini: deret pangkat
(\cref{ch:b2:powerseries}) adalah \omterm{def:b2:series:summable}{keluarga terjumlahkan} yang menyamar,
koefisien Fourier dikalikan lewat hasil kali Cauchy dan ditata
ulang oleh Parseval (\cref{ch:b2:fourier}), dan fungsi pembangkit
peluang (\cref{ch:b2:genfun}) tak lain teorema Fubini yang
diterapkan pada nilai harapan. Jilid Tahun ke-3 menyerap \omterm{def:b2:series:summable}{keluarga
terjumlahkan} ke dalam pengintegralan Lebesgue atas ukuran pencacah ---
tempat \cref{thm:b2:series:fubini} menjadi kasus khusus
teorema Fubini--Tonelli.
\end{remark}

\section{Latihan}

\begin{exercise}[$\star$]\label{exo:b2:series:1}
Sifat deret: $\sum \dfrac{\cos n}{n}$; $\;\sum \dfrac{(-1)^n}{\ln
n}$; $\;\sum \dfrac{(-1)^n}{n^{3/4} + \cos n}$ \emph{(uraikan seperti
pada jebakan Tahun ke-1: uji berselang-seling menuntut kemonotonan)}.
\end{exercise}

\begin{exercise}[$\star$]\label{exo:b2:series:2}
Buktikan bahwa untuk $\abs z < 1$: $\sum_{n\geq1} n z^{n} =
\dfrac{z}{(1-z)^2}$, lewat hasil kali Cauchy $\sum z^n$ dengan
dirinya sendiri.
\end{exercise}

\begin{exercise}[$\star\star$]\label{exo:b2:series:3}
(Lema bergaya Kronecker) Misalkan $\sum b_n$ deret real yang konvergen.
Buktikan, lewat \omterm{thm:b2:series:abel}{penjumlahan Abel}, bahwa $\dfrac{1}{n}\sum_{k=1}^{n}
k\,b_k \to 0$.
\end{exercise}

\begin{exercise}[$\star\star$]\label{exo:b2:series:4}
Telaahlah kekonvergenan $\sum \dfrac{\sin(n\theta)}{n^\alpha}$
($\theta \in \R$, $\alpha > 0$) --- untuk $(\theta, \alpha)$ yang mana ia
konvergen \omterm{def:b2:series:def}{mutlak}, semi-konvergen, atau divergen?
\end{exercise}

\begin{exercise}[$\star\star\star$]\label{exo:b2:series:5}
(Penataan ulang Riemann) Misalkan $\sum u_n$ deret real yang konvergen
tetapi tidak konvergen \omterm{def:b2:series:def}{mutlak}, dan $\ell \in \R$. Buktikan bahwa
suatu penataan ulang $\sum u_n$ konvergen ke $\ell$.
\emph{(Tunjukkan kedua deret bagian berisi suku positif dan suku negatif
sama-sama divergen; lalu berselang-selinglah secara rakus: ambil suku
positif sampai melampaui $\ell$, lalu suku negatif sampai turun di
bawahnya, dan seterusnya; sukunya menuju $0$, sehingga
kekonvergenannya ke $\ell$ terpaksa terjadi.)}
\end{exercise}

\begin{exercise}[$\star\star$]\label{exo:b2:series:6}
Buktikan bahwa keluarga
$\Bigl(\dfrac{x^{m+n}}{m!\,n!}\Bigr)_{(m,n)\in\N^2}$ \omterm{def:b2:series:summable}{terjumlahkan}
untuk setiap $x \in \R$, lalu turunkan kembali kesamaan $(\eu^x)^2 =
\eu^{2x}$ dengan mengelompokkan jumlah gandanya sepanjang diagonal $m +
n = k$.
\end{exercise}

\begin{exercise}[$\star\star$]\label{exo:b2:series:7}
Buktikan bahwa keluarga $\bigl(\frac{1}{m^2 n^2}\bigr)_{m,n \geq 1}$
\omterm{def:b2:series:summable}{terjumlahkan}, dan bahwa pengelompokan menurut $\gcd$, dengan $q = \gcd(m,n)$,
\[
\zeta(2)^2 = \sum_{q\geq1} \frac{1}{q^4}
\sum_{\substack{a,b \geq 1\\ \gcd(a,b)=1}} \frac{1}{a^2b^2}
= \zeta(4) \cdot S,
\]
dengan $S = \sum_{\gcd(a,b)=1} \frac{1}{a^2b^2}$: lalu turunkan $S =
\zeta(2)^2/\zeta(4)$. \emph{(Setiap pasangan $(m,n)$ tertulis secara
tunggal sebagai $(qa, qb)$ dengan $\gcd(a,b) = 1$.)}
\end{exercise}

\begin{exercise}[$\star\star\star$]\label{exo:b2:series:8}
(Teorema Abel tentang hasil kali, versi ringan) Andaikan $\sum a_n$
konvergen \omterm{def:b2:series:def}{mutlak} dan $\sum b_n$ konvergen. Buktikan bahwa hasil kali
Cauchy keduanya $\sum c_n$ konvergen, dengan $\sum c_n = (\sum
a_n)(\sum b_n)$. \emph{(Tulis $C_N = \sum_{k\leq N} c_k = \sum_n
a_n B_{N-n}$ dengan $B$ menyatakan jumlah parsial $b$; lalu pecah menurut
$n \leq N/2$ atau tidak, dengan memakai keterbatasan $(B_m)$ dan ekor
\omterm{def:b2:series:def}{mutlak} $(a_n)$.)}
\end{exercise}

\begin{exercise}[$\star\star\star$]\label{exo:b2:series:9}
Pada ruang $E$ (yang tidak \omterm{def:b2:metric:complete}{lengkap}) berisi barisan real yang akhirnya
nol dengan \omterm{def:b2:nvs:norm}{norma} supremum, tunjukkanlah sebuah deret yang konvergen
\omterm{def:b2:series:def}{mutlak} tetapi tidak konvergen di $E$. \emph{(Cobalah $u_n = 2^{-n} e_n$
dengan $(e_n)$ menyatakan barisan kanoniknya.)}
\end{exercise}

\begin{exercise}[$\star\star$]\label{exo:b2:series:10}
Periksalah kesamaan $\arctan(n+1) - \arctan(n) =
\arctan\dfrac{1}{n^2 + n + 1}$, lalu turunkan nilai eksak
\[
\sum_{n=1}^{\infty} \arctan\frac{1}{n^2 + n + 1} .
\]
\end{exercise}

\begin{exercise}[$\star\star$]\label{exo:b2:series:11}
Tentukan sifatnya (beserta yang setara) bagi
\[
\sum_n \bigl(\sqrt{n+1} - \sqrt n\bigr)^{\alpha}
\ (\alpha > 0), \qquad
\sum_n \Bigl(1 - \cos\frac1n\Bigr), \qquad
\sum_n \Bigl(\eu - \Bigl(1 + \frac1n\Bigr)^{\!n}\Bigr).
\]
\end{exercise}

\begin{exercise}[$\star\star\star$]\label{exo:b2:series:12}
Misalkan $(a_n)$ positif dan \emph{turun} dengan $\sum a_n$
konvergen. Buktikan bahwa $n\,a_n \to 0$ \emph{(batasi $n a_{2n}$ oleh
sebuah ekor)}. Tunjukkan bahwa konversnya gagal, dan bahwa
kemonotonannya penting, beserta contoh penyangkal yang gamblang.
\end{exercise}

\section{Soal: $\zeta(2) = \pi^2/6$ milik Euler, lewat Jumlah
Kotangen Cauchy}

Kesamaan Euler yang paling masyhur, $1 + \frac14 + \frac19 + \cdots =
\frac{\pi^2}6$, menerima bukti yang sepenuhnya dasar, karya
Cauchy: rumus de Moivre menghasilkan polinomial yang akarnya adalah
bilangan $\cot^2\frac{k\pi}{2n+1}$, Vieta menjumlahkan akar itu
secara eksak, lalu apitan $\cot^2\theta < \frac{1}{\theta^2} <
1 + \cot^2\theta$ meremukkan jumlah parsial $\sum\frac1{k^2}$
di antara dua batas rasional yang gamblang. Kita jalankan buktinya
selengkapnya, lalu sarikan $\zeta(4) = \frac{\pi^4}{90}$ dengan metode
yang sama, lalu petakan seluruh perbatasan antara kekonvergenan dan
kedivergenan lewat deret Bertrand --- dan buktikan bahwa perbatasan itu
sama sekali tak mengusung deret konvergen yang paling lambat.

\begin{problem}[{Soal akhir pekan --- $\zeta(2) = \pi^2/6$ dan
panorama Bertrand}]
\label{pb:b2:series:1}
Di sepanjang soal ini, $n \geq 1$ dan $\theta_k = \dfrac{k\pi}{2n+1}$
untuk $k = 1, \dots, n$; perhatikan $0 < \theta_k < \frac\pi2$.

\medskip
\noindent\textbf{Bagian I --- Kesamaan kotangen.}
\begin{enumerate}
  \item Buktikan rumus de Moivre $(\cos\theta +
        \iu\sin\theta)^m = \cos m\theta + \iu\sin m\theta$ ($m
        \in \N$), lalu turunkan, untuk $m = 2n + 1$,
        \[
        \sin\bigl((2n{+}1)\theta\bigr) = \sum_{j=0}^{n}
        (-1)^j\binom{2n+1}{2j+1}\cos^{2(n-j)}\theta\,
        \sin^{2j+1}\theta .
        \]
  \item Turunkan bahwa untuk $\theta \in \intoo{0}{\frac\pi2}$,
        \[
        \sin\bigl((2n{+}1)\theta\bigr) =
        \sin^{2n+1}\theta\; P_n(\cot^2\theta),
        \qquad
        P_n(x) = \sum_{j=0}^{n}
        (-1)^j\binom{2n+1}{2j+1}x^{\,n-j},
        \]
        yakni polinomial berderajat $n$ dengan koefisien utama $2n
        + 1$.
  \item Tunjukkan bahwa $x_k = \cot^2\theta_k$, $k = 1, \dots, n$,
        merupakan $n$ akar $P_n$ yang \emph{berbeda} --- jadi
        seluruh akarnya.
  \item Lewat Vieta, buktikan kesamaan eksak
        \[
        \sum_{k=1}^{n} \cot^2\frac{k\pi}{2n+1}
        = \frac{n(2n-1)}{3}.
        \]
  \item Turunkan pula $\displaystyle\sum_{k=1}^{n}
        \frac{1}{\sin^2\theta_k} = \frac{2n(n+1)}{3}$.
  \item Buktikan apitannya: $\cot^2\theta < \dfrac1{\theta^2} <
        \dfrac{1}{\sin^2\theta}$ untuk $\theta \in
        \intoo{0}{\frac\pi2}$ \emph{(dari $\sin\theta < \theta <
        \tan\theta$)}.
\end{enumerate}

\medskip
\noindent\textbf{Bagian II --- Apitannya menutup: teorema
Euler.}
\begin{enumerate}[resume]
  \item Dengan menjumlahkan pertanyaan 6 atas $k = 1, \dots, n$ pada
        $\theta = \theta_k$, tegakkan
        \[
        \frac{n(2n-1)}{3} \;<\;
        \frac{(2n+1)^2}{\pi^2}\sum_{k=1}^{n}\frac1{k^2}
        \;<\; \frac{2n(n+1)}{3}.
        \]
  \item Simpulkan (\emph{teorema Euler, lewat bukti Cauchy}):
        \[
        \zeta(2) = \sum_{k=1}^{\infty}\frac{1}{k^2} =
        \frac{\pi^2}{6}.
        \]
  \item Sarikan lajunya dari apitan itu: tunjukkan
        \[
        \Bigl|\sum_{k=1}^{n}\frac1{k^2} - \frac{\pi^2}6\Bigr|
        = O\Bigl(\frac1n\Bigr),
        \]
        yang selaras dengan ekor eksak $\sum_{k>n}k^{-2} =
        \frac1n - \frac{1}{2n^2} + O(n^{-3})$ pada
        \cref{exo:b2:comparison:11}.
  \item Jalankan mesinnya satu lantai lebih tinggi: dengan memakai
        fungsi Vieta kedua bagi $P_n$, tunjukkan
        \[
        \sum_{k=1}^{n}\cot^4\theta_k =
        \Bigl(\frac{n(2n-1)}3\Bigr)^{\!2} -
        \frac{2n(2n-1)(2n-2)(2n-3)}{60}
        \;\sim\; \frac{8n^4}{45},
        \]
        lalu apitlah dengan $\cot^4 < \theta^{-4} < (1 +
        \cot^2)^2$ untuk memperoleh $\zeta(4) = \dfrac{\pi^4}{90}$.
\end{enumerate}

\medskip
\noindent\textbf{Bagian III --- Panennya.}
\begin{enumerate}[resume]
  \item Turunkan dari $\zeta(2) = \frac{\pi^2}6$:
        \[
        \sum_{k\geq0}\frac{1}{(2k+1)^2} = \frac{\pi^2}{8},
        \qquad
        \sum_{k\geq1}\frac{(-1)^{k-1}}{k^2} = \frac{\pi^2}{12}.
        \]
  \item Padukan dengan \cref{exo:b2:series:7}: hitunglah $S =
        \sum_{\gcd(a,b)=1}\frac{1}{a^2b^2} =
        \frac{\zeta(2)^2}{\zeta(4)} = \frac52$, lalu tafsirkan
        $\frac{1}{\zeta(2)} = \frac{6}{\pi^2} \approx 0.608$ sebagai
        kerapatan pasangan yang saling prima (nyatakan heuristiknya
        secara jujur: pencacahan yang cermat adalah urusan jilid Tahun
        ke-3).
  \item (Percepatan yang tersahkan) Rumus ekor pada pertanyaan 9
        memberi $\sum_{k\leq n}k^{-2} + \frac1n - \frac1{2n^2} =
        \frac{\pi^2}6 + O(n^{-3})$. Bandingkan kerja yang diperlukan
        untuk enam angka $\zeta(2)$: penjumlahan langsung lawan
        jumlah terkoreksi di $n = 100$ (tempat galatnya
        $1.7\cdot10^{-7}$).
  \item Periksalah pertanyaan 4 dengan tangan di $n = 1$ dan $n = 2$
        (yakni nilai $\cot^2\frac\pi3 = \frac13$ dan
        $\cot^2\frac\pi5 + \cot^2\frac{2\pi}5 = 2$), dengan memakai
        $\cos\frac\pi5 = \frac{1+\sqrt5}4$ atau penilaian numerik.
  \item Buktikan kesamaan pendampingnya
        \[
        \sum_{k=1}^{n}\tan^2\frac{k\pi}{2n+1} = n(2n+1)
        \]
        \emph{(sebab bilangan $\tan^2\theta_k$ adalah akar
        polinomial terbalik $x^nP_n(1/x)$)}, lalu periksalah di $n
        = 1$.
\end{enumerate}

\medskip
\noindent\textbf{Bagian IV --- Panorama Bertrand.} Untuk $\alpha,
\beta \in \R$, tinjaulah \emph{deret Bertrand}
\[
\sum_{n \geq 3} \frac{1}{n^{\alpha}(\ln n)^{\beta}} .
\]
\begin{enumerate}[resume]
  \item Tunjukkan bahwa untuk $\alpha > 1$ deretnya konvergen,
        berapa pun $\beta$ \emph{(bandingkan dengan
        $n^{-(1+\alpha)/2}$)}.
  \item Tunjukkan bahwa untuk $\alpha < 1$ ia divergen, berapa pun
        $\beta$.
  \item Untuk $\alpha = 1$: dengan memakai perbandingan deret dengan
        integral (\cref{thm:b2:comparison:seriesintegral}) pada $f(t) =
        \frac{1}{t(\ln t)^\beta}$, buktikan bahwa ia konvergen bila dan
        hanya bila $\beta > 1$.
  \item Iterasikan perbatasannya: tunjukkan $\sum\frac{1}{n\ln n\,\ln\ln
        n}$ divergen sedangkan $\sum\frac{1}{n\ln n\,(\ln\ln n)^2}$
        konvergen.
  \item Dua jebakan: tentukan sifat
        \[
        \sum_n \frac{1}{n^{1 + 1/\ln n}}
        \qquad\text{dan}\qquad
        \sum_n \frac{1}{n^{1 + 1/\ln\ln n}}
        \]
        \emph{(hitunglah $n^{1/\ln n}$ secara eksak; lalu bandingkan
        $n^{1/\ln\ln n}$ dengan setiap pangkat $\ln n$)}.
  \item (Tak ada deret konvergen yang paling lambat) Misalkan $\sum a_n$
        deret konvergen apa pun dengan $a_n > 0$, dan $R_n = \sum_{k
        \geq n}a_k$ ekornya. Buktikan bahwa $\sum
        \frac{a_n}{\sqrt{R_n}}$ tetap \emph{konvergen}
        \emph{(bandingkan dengan teleskop $2(\sqrt{R_n} -
        \sqrt{R_{n+1}})$)}, walaupun $\frac{a_n/\sqrt{R_n}}{a_n}
        \to \infty$: jadi setiap deret konvergen didominasi tegas
        oleh deret konvergen yang lain. Perbatasan
        kekonvergenannya bukan sebuah kurva, melainkan kabut.
\end{enumerate}

\medskip
\noindent\textbf{Bagian V --- Pemeriksaan silang dan rangkuman.}
\begin{enumerate}[resume]
  \item (Kondensasi Cauchy) Buktikan: untuk $(a_n)$ yang positif dan
        turun, $\sum a_n$ konvergen bila dan hanya bila $\sum 2^k a_{2^k}$
        konvergen. Lalu turunkan kembali perbatasan pertanyaan 18 dari situ.
  \item (Ongkos kelambanan) Untuk $\sum\frac1{n(\ln n)^2}$,
        batasilah ekornya oleh sebuah integral lalu tunjukkan bahwa
        menjumlahkannya sampai $N = 10^6$ masih menyisakan galat yang
        lebih besar daripada $0.07$: jadi kekonvergenan yang disahkan
        teori dapat tak berguna bagi numerik --- bandingkan dengan
        pertanyaan 13.
  \item Golongkanlah (beserta pembenaran satu baris):
        $\sum\frac1{n\ln n}$, $\sum\frac1{n^{1.01}}$,
        $\sum\frac{(\ln n)^{100}}{n^{1.001}}$, dan
        $\sum\frac1{n(\ln n)(\ln\ln n)^{3}}$.
  \item (Rangkuman) Satu kalimat untuk masing-masing: bagaimana de
        Moivre mengubah kesamaan trigonometri menjadi polinomial yang
        jumlah akarnya terhitungkan; di mana apitannya memerlukan
        kesamaan \emph{eksak} di ujungnya, bukan sekadar yang setara;
        perkakas \cref{ch:b2:comparison} yang mana yang menggerakkan
        Bagian IV; dan apa yang dikatakan pertanyaan 21 tentang impian
        adanya ``uji perbandingan semesta''. Sebutkan kedua puncaknya:
        $\zeta(2) = \frac{\pi^2}6$ milik Euler (beserta lantai di
        atasnya, $\zeta(4) = \frac{\pi^4}{90}$), dan penggolongan
        Bertrand. Catat pula di mana $\zeta(2)$ akan dibuktikan
        sekali lagi: lewat Parseval pada \cref{ch:b2:fourier} --- satu
        teorema, dua peradaban.

\end{enumerate}
\end{problem}
